\exnumber{00}
\extitle{UI III - Ex 1 - Grilles}
\exfiles{ex00/fichier Figma, nomme Name\_FirstName\_DeliverableName\_Date\_VersionNumber}
\exauthorize{Ceux nommés dans l'exercice}
\makeheaderfiles

\textbf{Contexte :}

Bienvenue dans la 3e quête de UI.

Dans cette série d'exercices, on vous demandera d'ouvrir les yeux et de collecter des éléments graphiques sur divers sites web. Vous verrez des grilles à travers les écrans, des boutons multi-états, comment rendre un design accessible pour tous, des fils d'Ariane, des boutons toggle, des boutons radio, un calendrier, un time picker, avec 2 objectifs principaux :

\begin{itemize}
\item Enrichir votre culture générale graphique
\item Détecter la cohérence, qui aide les utilisateurs à comprendre un produit numérique dès le premier regard.
\end{itemize}

\textbf{Instructions :}

Ouvrez ces 5 images sur Figma et tracez des lignes et des grilles chaque fois que vous trouvez des éléments alignés.

\textbf{Livrables :}

\begin{itemize}
\item Un fichier Figma comprenant les écrans avec des lignes par-dessus pour montrer les grilles. Une page par écran ou frame.
\item Déposez le fichier Figma sur Github, nommé ainsi : « Name\_FirstName\_DeliverableName\_Date\_VersionNumber ». Exemple : « Doe\_John\_MidFiPrototype\_05242024\_V1 ».
\item Ajoutez un titre dans le fichier.
\item Soignez l'aspect global du fichier. Il doit être clair, simple et lisible. Vous pouvez regarder canva.com, sans surcharger le dessin de détails.
\end{itemize}

\textbf{Ressources :}

\begin{itemize}
\item Tutoriel vidéo sur Figma \url{https://www.youtube.com/watch?v=FTFaQWZBqQ8}
\item Débuter en design graphique \url{https://edu.gcfglobal.org/en/beginning-graphic-design/}
\item Gareth graphic studio \url{https://www.youtube.com/c/GarethDavidStudio}
\item Exemple du 1er écran \url{https://www.figma.com/file/KOzd8yZuJvGLmZLuXjBgDg/UI-III-Ex-1.1}
\end{itemize}

